%% kapa: exam=1 round=31 year=2020 period=2 form=A subject=appraisal_law subject_ko=감정평가관계법규 seq=27 number=27
%% source: Q-Net 감정평가사 기출문제 (cst00309) artlSeq=5207219 file=extracted/1008/5207219/A형/A형/2020년도 제31회 감정평가사 1차 시험_2교시_A형.pdf
%% conversion: kordoc 4.24.51 -> tools/gen1.py (auto); review=proofread-sonnet; flags=none
\begin{question}{27}
건축법령상 건축물의 구조 및 재료 등에 관한 설명으로 옳지 않은 것은?
\begin{choices}
  \item 건축물은 고정하중, 적재하중, 적설하중, 풍압, 지진, 그 밖의 진동 및 충격 등에 대하여 안전한 구조를 가져야 한다.
  \item 지방자치단체의 장은 구조 안전 확인 대상 건축물에 대하여 건축허가를 하는 경우 내진성능 확보 여부를 확인하여야 한다.
  \item 국토교통부장관은 지진으로부터 건축물의 구조 안전을 확보하기 위하여 건축물의 용도, 규모 및 설계구조의 중요도에 따라 내진등급을 설정하여야 한다.
  \item 연면적이 200제곱미터인 목구조 건축물을 건축하고자 하는 자는 사용승인을 받는 즉시 내진능력을 공개하여야 한다.
  \item 국가 또는 지방자치단체는 건축물의 소유자나 관리자에게 피난시설 등의 설치, 개량ㆍ보수 등 유지ㆍ관리에 대한 기술지원을 할 수 있다.
\end{choices}
\end{question}
\answer{4}
